\documentclass[12pt]{article}
\usepackage[margin=2.5cm]{geometry}
\usepackage{amsmath, amssymb}
\usepackage{bm}
\usepackage{parskip}

\title{Likelihood for ascertained proband--relative--unrelated strata}
\author{}
\date{}

\begin{document}
\maketitle

\section*{Sampling design}

Families containing at least one case are ascertained.  From each ascertained
family we observe three individuals:
\begin{itemize}
  \item Proband $P$: the index case, so $y_P = 1$ by design.
  \item Relative $R$: one family member of $P$, observed with their natural
        disease status $y_R \in \{0,1\}$.
  \item Unrelated $U$: one individual sampled at random from the population,
        observed with their natural disease status $y_U \in \{0,1\}$.
\end{itemize}
Let $G_P$, $G_R$, $G_U$ denote the unobserved genotypic values of these three
individuals.

\section*{Logistic model}

Under the logistic model, outcomes are conditionally independent given genotypic
values:
\[
  P(y_i = 1 \mid G_i) = \sigma(\beta_0 + G_i)
  \;=\; \frac{e^{\,\beta_0 + G_i}}{1 + e^{\,\beta_0 + G_i}},
\]
where $\beta_0$ is the population log-odds of disease.  The genotypic values
satisfy
\[
  G_i \;=\; s\,(L\bm{z})_i, \qquad \bm{z} \sim \mathcal{N}(\bm{0},\bm{I}),
\]
with $L$ the Cholesky factor of the genetic relationship matrix and
$\mu = \tfrac{1}{2}s^2$ the genetic information parameter.

\section*{Joint likelihood conditional on genotypic values}

Given $G_P, G_R, G_U$, the three outcomes are independent:
\begin{equation}
  P(y_P, y_R, y_U \mid G_P, G_R, G_U)
  \;=\;
  \sigma(\beta_0 + G_P)^{y_P}[1-\sigma(\beta_0+G_P)]^{1-y_P}
  \cdot
  P(y_R \mid G_R)
  \cdot
  P(y_U \mid G_U),
\end{equation}
with the proband factor evaluated at $y_P = 1$.

\section*{Ascertainment correction}

Because families are selected on having a case, the ascertainment probability
for a family containing at least one member with genotypic value $G_P$ is
\[
  P(\text{ascertained} \mid G_P) \;=\; P(y_P = 1 \mid G_P)
  \;=\; \sigma(\beta_0 + G_P).
\]
The ascertainment-corrected likelihood divides the joint probability by this
selection probability:
\begin{align}
  L(\beta_0, s \mid y_R, y_U,\, G_P, G_R, G_U)
  &\;\propto\;
  \frac{P(y_P=1, y_R, y_U \mid G_P, G_R, G_U)}
       {P(\text{ascertained} \mid G_P)}
  \notag\\[6pt]
  &\;=\;
  \frac{\sigma(\beta_0+G_P)\cdot P(y_R \mid G_R)\cdot P(y_U \mid G_U)}
       {\sigma(\beta_0+G_P)}.
\end{align}

The proband factor $\sigma(\beta_0+G_P)$ cancels, giving
\begin{equation}
  \boxed{
    L(\beta_0, s \mid y_R, y_U,\, G_P, G_R, G_U)
    \;=\;
    P(y_R \mid G_R)\cdot P(y_U \mid G_U).
  }
\end{equation}

The proband's genotypic value $G_P$ does \emph{not} appear in the likelihood.
It enters inference only through the prior on the genotypic values: because $P$
and $R$ are relatives, $(G_P, G_R)$ are jointly Gaussian with covariance
$s^2 \cdot r_{PR}$, where $r_{PR}$ is the kinship coefficient (e.g.\
$r_{PR}=0.5$ for full siblings).  Conditioning on $y_P = 1$ in the joint
prior updates the distribution of $G_P$ upward, which in turn shifts the prior
on $G_R$ upward:
\[
  G_R \mid G_P \;\sim\; \mathcal{N}\!\left(r_{PR}\,G_P,\;
  s^2(1 - r_{PR}^2)\right).
\]

\section*{Why the proband's case status informs the relative's risk}

Marginalising over $G_P$ and $G_R$ (i.e.\ integrating over the genotypic
prior):
\[
  P(y_R = 1 \mid y_P = 1)
  \;=\;
  \int P(y_R=1 \mid G_R)\;
  p(G_R \mid G_P)\;
  p(G_P \mid y_P=1)\;
  dG_P\, dG_R
  \;>\; K,
\]
where $K = P(y_R=1)$ is the population prevalence.  Observing $y_P=1$
shifts the posterior on $G_P$ toward higher values (ascertainment), which
shifts the prior on $G_R$ upward, which elevates $P(y_R=1 \mid y_P=1)$.
The ratio
\[
  \lambda_S \;=\; \frac{P(y_R=1 \mid y_P=1)}{K}
  \;\approx\; e^{\mu} \quad \text{for small } K \text{ and } \mu,
\]
is the sibling recurrence risk ratio.  The proband's case status is thus
informative about the relative's risk only through this shared-genetics
channel, not through the ascertainment-corrected likelihood.

\section*{Expanded form of the likelihood}

Writing out the two logistic factors:
\begin{align}
  L
  &\;=\;
  \sigma(\beta_0 + G_R)^{y_R}
  \left[1 - \sigma(\beta_0 + G_R)\right]^{1-y_R}
  \cdot
  \sigma(\beta_0 + G_U)^{y_U}
  \left[1 - \sigma(\beta_0 + G_U)\right]^{1-y_U}.
\end{align}

Both $\beta_0$ and $s$ (through its effect on $G_R$ and $G_U$) are
estimable:
\begin{itemize}
  \item \textbf{Intercept $\beta_0$}: identified by the overall case rate among
        $R$ and $U$.
  \item \textbf{Scale $s$ (genetic signal)}: identified by the excess case rate
        in $R$ relative to $U$, given that the proband was a case.
\end{itemize}

\section*{Comparison with the conditional logistic likelihood}

If instead we ignored the sampling design and treated $(P, R, U)$ as a
matched set of three, the conditional logistic likelihood would condition on
$\sum_i y_i = 1 + y_R + y_U$ and give equal weight to all permutations of
case assignment within the stratum (see \textit{cond\_likelihood.pdf}).  This
is incorrect here because the proband's case status is \emph{fixed by design},
not interchangeable with the other stratum members.  The ascertainment-corrected
likelihood derived above accounts for the asymmetric roles of $P$, $R$, and $U$.

\end{document}
